\documentclass[10pt,a4paper]{article}

\usepackage[a4paper,top=13mm,bottom=14mm,left=14mm,right=14mm]{geometry}
\usepackage{fontspec}
\usepackage{xcolor}
\usepackage{tabularx}
\usepackage{array}
\usepackage[most]{tcolorbox}
\usepackage{fancyhdr}
\usepackage{lastpage}
\usepackage{enumitem}
\usepackage{amssymb}
\usepackage{microtype}

\IfFontExistsTF{Arial}{
  \setmainfont{Arial}
  \setsansfont{Arial}
}{
  \setmainfont{TeX Gyre Heros}
  \setsansfont{TeX Gyre Heros}
}

\definecolor{AIEPRed}{HTML}{D62027}
\definecolor{AIEPNavy}{HTML}{082B5C}
\definecolor{AIEPInk}{HTML}{182B3A}
\definecolor{AIEPMuted}{HTML}{5F6B7A}
\definecolor{AIEPLine}{HTML}{C9D2DE}
\definecolor{AIEPSoft}{HTML}{F5F2EC}
\definecolor{AIEPBlueSoft}{HTML}{E9EEF4}
\definecolor{AIEPCyan}{HTML}{35B7C6}
\definecolor{AIEPGold}{HTML}{E0BC5A}
\definecolor{AIEPGreen}{HTML}{2E8B57}

\pagestyle{fancy}
\fancyhf{}
\renewcommand{\headrulewidth}{0pt}
\fancyfoot[L]{\footnotesize\textcolor{AIEPMuted}{GeoGreen Escolar · Mentoría 2}}
\fancyfoot[R]{\footnotesize\textcolor{AIEPMuted}{Ficha de trabajo por equipo · \thepage/\pageref{LastPage}}}

\setlength{\parindent}{0pt}
\setlength{\parskip}{3pt}
\renewcommand{\arraystretch}{1.18}
\setlist[itemize]{leftmargin=*,topsep=1pt,itemsep=1pt}
\newcolumntype{Y}{>{\raggedright\arraybackslash}X}
\newcolumntype{C}[1]{>{\centering\arraybackslash}p{#1}}

\newtcolorbox{fieldbox}[2][]{
  colback=white,
  colframe=AIEPLine,
  boxrule=0.55pt,
  arc=1mm,
  left=2.5mm,right=2.5mm,top=2mm,bottom=2mm,
  before skip=2pt,after skip=2pt,
  title={\scriptsize\bfseries\MakeUppercase{#2}},
  coltitle=AIEPNavy,
  fonttitle=\sffamily,
  attach title to upper={\par\vspace{0.8mm}},
  #1
}

\newtcolorbox{accentfield}[2][]{
  colback=AIEPSoft,
  colframe=AIEPRed,
  boxrule=0.7pt,
  arc=1mm,
  left=2.5mm,right=2.5mm,top=2mm,bottom=2mm,
  before skip=2pt,after skip=2pt,
  title={\scriptsize\bfseries\MakeUppercase{#2}},
  coltitle=AIEPRed,
  fonttitle=\sffamily,
  attach title to upper={\par\vspace{0.8mm}},
  #1
}

\newcommand{\writearea}[1]{\vspace{#1}\hrule height 0.35pt\vspace{4mm}\hrule height 0.35pt}
\newcommand{\checkline}[1]{$\square$\hspace{1.5mm}#1}
\newcommand{\pagehead}[2]{%
  {\small\bfseries\textcolor{AIEPRed}{GEOGREEN ESCOLAR · MENTORÍA 2}}\par
  \vspace{2mm}
  {\LARGE\bfseries\color{AIEPNavy}#1\par}
  {\normalsize\color{AIEPMuted}#2\par}
  \vspace{4mm}
}

\begin{document}

\pagehead{Definir el producto mínimo viable}{Solución propuesta y recursos definidos · Conservar para la Mentoría 3}

\begin{tabularx}{\textwidth}{@{}>{\bfseries\color{AIEPNavy}}p{25mm}Y >{\bfseries\color{AIEPNavy}}p{29mm}p{38mm}@{}}
Equipo & \rule{57mm}{0.35pt} & Proyecto & \rule{36mm}{0.35pt} \\
Integrantes & \multicolumn{3}{l}{\rule{145mm}{0.35pt}} \\
\end{tabularx}

\vspace{3mm}
\begin{tcolorbox}[colback=AIEPNavy,colframe=AIEPNavy,coltext=white,boxrule=0pt,arc=1mm,left=3mm,right=3mm,top=2mm,bottom=2mm]
\textbf{Productos de entrada:}\hspace{4mm}
\checkline{Problema validado M1}\hspace{5mm}
\checkline{Idea tecnológica T3}\hspace{5mm}
\checkline{Próximo paso}
\end{tcolorbox}

{\large\bfseries\color{AIEPNavy}1. Del problema a la solución}\par

\begin{fieldbox}[height=31mm]{Problema validado y contexto afectado}
¿Qué ocurre, dónde ocurre y quién necesita una respuesta o información?
\writearea{5mm}
\end{fieldbox}

\begin{accentfield}[height=40mm]{Nuestra solución en una frase}
Cuando \rule{37mm}{0.35pt} pase de \rule{29mm}{0.35pt}, entonces \rule{43mm}{0.35pt},\\[5mm]
para que \rule{42mm}{0.35pt} pueda \rule{66mm}{0.35pt}.
\end{accentfield}

{\large\bfseries\color{AIEPNavy}2. Regla del sistema}\par
\begin{tabularx}{\textwidth}{@{}Y@{\hspace{4mm}}Y@{}}
\begin{fieldbox}[height=35mm]{Variable y unidad}
¿Qué cambia y cómo se mide o distingue?
\writearea{7mm}
\end{fieldbox}
&
\begin{fieldbox}[height=35mm]{Sensor o fuente de datos}
¿Con qué se obtiene la lectura y por qué sirve?
\writearea{7mm}
\end{fieldbox}
\\
\begin{fieldbox}[height=35mm]{Umbral y origen}
¿Cuándo se activa? Marquen:\par
\checkline{medición}\quad\checkline{cálculo}\quad\checkline{fuente}\quad\checkline{provisional}
\writearea{5mm}
\end{fieldbox}
&
\begin{fieldbox}[height=35mm]{Respuesta perceptible}
¿Qué luz, sonido, texto o acción cambia?
\writearea{7mm}
\end{fieldbox}
\end{tabularx}

\begin{fieldbox}[height=34mm]{Mostrar: destinatario, información y decisión}
¿Dónde se verá el resultado, quién lo mirará y qué podrá decidir con esa información?
\writearea{6mm}
\end{fieldbox}

\begin{tcolorbox}[colback=AIEPBlueSoft,colframe=AIEPLine,boxrule=0.5pt,arc=1mm,left=3mm,right=3mm,top=2mm,bottom=2mm]
\textbf{Comprobación:} \checkline{Mide algo real}\quad
\checkline{Decide con una regla}\quad
\checkline{Responde}\quad
\checkline{Se muestra a otra persona}
\end{tcolorbox}

\newpage
\pagehead{Recursos, software y seguridad}{Definir lo disponible, lo que debe conseguirse y lo que se representará}

{\large\bfseries\color{AIEPNavy}3. Lista de recursos con función}\par
{\small Cada recurso debe responder a una parte de la regla. Incluyan componentes, materiales, herramientas, permisos, espacios y conocimientos.}

\small
\begin{tabularx}{\textwidth}{|p{38mm}|p{49mm}|Y|}
\hline
\textbf{Ya lo tenemos} & \textbf{Debemos conseguirlo}\newline Recurso · responsable · fecha & \textbf{Lo representaremos}\newline Simulación · maqueta · diagrama · registro \\
\hline
\rule{0pt}{18mm} & & \\
\hline
\rule{0pt}{18mm} & & \\
\hline
\rule{0pt}{18mm} & & \\
\hline
\end{tabularx}
\normalsize

\begin{fieldbox}[height=26mm]{Función de los recursos principales}
¿Qué parte de observar, decidir, responder o mostrar cumple cada recurso?
\writearea{4mm}
\end{fieldbox}

{\large\bfseries\color{AIEPNavy}4. Capa de software}\par
\begin{tabularx}{\textwidth}{@{}Y@{\hspace{4mm}}Y@{}}
\begin{fieldbox}[height=28mm]{Datos que entran}
¿Qué valor o estado recibe y de dónde viene?
\writearea{5mm}
\end{fieldbox}
&
\begin{fieldbox}[height=28mm]{Qué debe mostrar}
¿Cuál es la información principal?
\writearea{5mm}
\end{fieldbox}
\\
\begin{fieldbox}[height=28mm]{Quién lo usa}
¿Qué persona lo mira y qué decide?
\writearea{5mm}
\end{fieldbox}
&
\begin{fieldbox}[height=28mm]{Dónde debe funcionar}
¿Pantalla, computador, navegador, planilla u otro?
\writearea{5mm}
\end{fieldbox}
\end{tabularx}

\begin{fieldbox}[height=25mm]{Apoyo utilizado y verificación del equipo}
¿Qué apoyo solicitaron, qué comprobaron y qué decisión tomaron?
\writearea{4mm}
\end{fieldbox}

{\large\bfseries\color{AIEPNavy}5. Revisión antes de energizar}\par
\begin{tcolorbox}[colback=AIEPNavy,colframe=AIEPNavy,coltext=white,boxrule=0pt,arc=1mm,left=3mm,right=3mm,top=2.5mm,bottom=2.5mm]
\checkline{Voltaje correcto}\hspace{5mm}
\checkline{Polaridad correcta}\hspace{5mm}
\checkline{Tierra común}\par
\checkline{Resistencias instaladas}\hspace{5mm}
\checkline{Montaje sin contactos indebidos}\hspace{5mm}
\checkline{Revisión cruzada}\par
\vspace{1.5mm}
\textbf{Si existe una duda, el montaje permanece sin energía y se consulta antes de conectar.}
\end{tcolorbox}

\newpage
\pagehead{Alcance y trabajo distribuido}{Una responsabilidad por integrante · Una tarea verificable por responsabilidad}

{\large\bfseries\color{AIEPNavy}6. Alcance comprometido}\par
\begin{tabularx}{\textwidth}{@{}Y@{\hspace{4mm}}Y@{}}
\begin{accentfield}[height=33mm]{Comprometido para el 21 de septiembre}
¿Qué versión mínima podremos mostrar?
\writearea{9mm}
\end{accentfield}
&
\begin{fieldbox}[height=33mm]{Fuera de alcance por ahora}
¿Qué función o mejora decidimos postergar?
\writearea{9mm}
\end{fieldbox}
\end{tabularx}

\begin{fieldbox}[height=28mm]{Plan alternativo}
Si falta un material, falla el sensor o no disponemos de la placa, ¿qué evidencia podremos mostrar igualmente?
\writearea{5mm}
\end{fieldbox}

{\large\bfseries\color{AIEPNavy}7. Plan de trabajo de las seis responsabilidades}\par
{\footnotesize Cada fila debe contener una acción concreta, una fecha anterior a la Mentoría 3 y una forma de comprobarla.}

\footnotesize
\begin{tabularx}{\textwidth}{|p{27mm}|p{28mm}|Y|p{20mm}|p{38mm}|}
\hline
\textbf{Responsabilidad} & \textbf{Integrante} & \textbf{Tarea concreta} & \textbf{Fecha} & \textbf{Cómo se comprueba} \\
\hline
\rule{0pt}{14mm}Coordinación & & & & \\
\hline
\rule{0pt}{14mm}Investigación & & & & \\
\hline
\rule{0pt}{14mm}Diseño & & & & \\
\hline
\rule{0pt}{14mm}Tecnología & & & & \\
\hline
\rule{0pt}{14mm}Pruebas y evidencia & & & & \\
\hline
\rule{0pt}{14mm}Comunicación & & & & \\
\hline
\end{tabularx}
\normalsize

\begin{tcolorbox}[colback=AIEPBlueSoft,colframe=AIEPLine,boxrule=0.5pt,arc=1mm,left=3mm,right=3mm,top=2mm,bottom=2mm]
\textbf{Orden de trabajo:}\quad
\checkline{Primero una lectura que reacciona}\quad
\checkline{Luego la regla}\quad
\checkline{Después la respuesta}\quad
\checkline{Finalmente integrar y registrar}
\end{tcolorbox}

\begin{accentfield}[height=26mm]{Turno con la placa o el equipo compartido}
Antes de nuestro turno dejaremos listo: \rule{99mm}{0.35pt}.\\[4mm]
Durante el turno probaremos únicamente: \rule{92mm}{0.35pt}.
\end{accentfield}

\newpage
\pagehead{Profundización, evidencia y continuidad}{La Mentoría 3 revisa lo construido, probado y registrado}

{\large\bfseries\color{AIEPNavy}8. Rama de profundización}\par
\begin{tcolorbox}[colback=AIEPSoft,colframe=AIEPLine,boxrule=0.5pt,arc=1mm,left=3mm,right=3mm,top=2mm,bottom=2mm]
\checkline{Simulación y verificación}\hspace{4mm}
\checkline{Electrónica y montaje}\hspace{4mm}
\checkline{Software y datos}\par
\checkline{Diseño físico y fabricación}\hspace{4mm}
\checkline{Comunicación y evidencia}
\end{tcolorbox}

\begin{fieldbox}[height=33mm]{Por qué esta rama fortalece nuestra propuesta}
¿Cómo se conecta con el problema y qué hará visible ante otra persona?
\writearea{5mm}
\end{fieldbox}

{\large\bfseries\color{AIEPNavy}9. Evidencia comprometida para la Mentoría 3}\par
\begin{tcolorbox}[colback=AIEPBlueSoft,colframe=AIEPLine,boxrule=0.5pt,arc=1mm,left=3mm,right=3mm,top=2mm,bottom=2mm]
\checkline{Prototipo parcial}\hspace{4mm}
\checkline{Simulación}\hspace{4mm}
\checkline{Diagrama verificado}\hspace{4mm}
\checkline{Maqueta}\par
\checkline{Registro de mediciones}\hspace{4mm}
\checkline{Panel o registro digital}\hspace{4mm}
\checkline{Otra: \rule{31mm}{0.35pt}}
\end{tcolorbox}

\begin{tabularx}{\textwidth}{@{}Y@{\hspace{4mm}}Y@{}}
\begin{fieldbox}[height=39mm]{Qué probaremos}
¿Qué cambio o funcionamiento debe observarse?
\writearea{9mm}
\end{fieldbox}
&
\begin{fieldbox}[height=39mm]{Qué registraremos}
¿Qué fotografía, video, dato, captura o resultado guardaremos?
\writearea{9mm}
\end{fieldbox}
\\
\begin{fieldbox}[height=39mm]{Qué esperamos aprender}
¿Qué decisión podremos mantener o corregir?
\writearea{9mm}
\end{fieldbox}
&
\begin{fieldbox}[height=39mm]{Qué mostraremos si falla}
¿Cómo convertiremos el error en evidencia útil?
\writearea{9mm}
\end{fieldbox}
\end{tabularx}

{\large\bfseries\color{AIEPNavy}10. Lista de salida obligatoria}\par
\begin{tcolorbox}[colback=AIEPNavy,colframe=AIEPNavy,coltext=white,boxrule=0pt,arc=1mm,left=3mm,right=3mm,top=2.5mm,bottom=2.5mm]
\checkline{Regla completa}\hspace{4mm}
\checkline{Recursos y seguridad}\hspace{4mm}
\checkline{Software y destinatario}\par
\checkline{Alcance y alternativa}\hspace{4mm}
\checkline{Seis tareas verificables}\hspace{4mm}
\checkline{Evidencia definida}
\end{tcolorbox}

\begin{accentfield}[height=42mm]{Declaración del equipo}
Nuestro producto mínimo viable va a \rule{91mm}{0.35pt}.\\[5mm]
Antes del 21 de septiembre vamos a demostrar que \rule{77mm}{0.35pt}.
\end{accentfield}

\begin{tcolorbox}[colback=AIEPSoft,colframe=AIEPRed,boxrule=0.7pt,arc=1mm,left=3mm,right=3mm,top=2mm,bottom=2mm]
\textbf{Conserven esta ficha.} Entre sesiones ejecuten las tareas, realicen pruebas y guarden evidencia. En la Mentoría 3 deberán mostrar qué probaron y qué aprendieron.
\end{tcolorbox}

\end{document}
